\chapter{Conclusions and Future Scope}
\label{Chapter 6}

\section{Conclusion}
This report has been the writeup of a five-month engineering internship at Ignosis. Over the course of those five months I contributed to a production, multi-tenant AI-Based Collections Platform. The platform is now live across four lender tenants and is materially improving the recovery numbers those lenders had historically been getting. Personally, the internship moved me from an undergraduate who could write Java and React but had never shipped anything serious, to someone now reasonably comfortable putting code into a mature multi-tenant backend, a real frontend dashboard, a Node.js automation toolkit and a desktop client. The other thing the internship taught me, which is harder to put on a resume, is the operational discipline of production engineering -- writing tests that hit the real database rather than mocks, writing plans before writing code, writing the alert before the incident, and designing for tenant isolation by default rather than as an afterthought.

The headline numbers sit in Table~\ref{tab:headline} of Chapter~\ref{Chapter 5}. Bucket~0 collection rate moved from 64 percent to 85 percent. Share of cases recovered inside the first five days of a cycle moved from 38 percent to 70 percent. The multi-channel cumulative contact rate landed at 96 percent. Those numbers are very obviously a team product, not an individual one -- but the modules I contributed (the case detail and case listing screens, the bulk campaign module, the speed dialer, the WhatsApp automation, the desktop client, the Android SMS sender) are direct enablers of each of those three figures. The project also met all nine of the engineering problems pulled out at the end of Chapter~\ref{Chapter 2}. None of the individual implementations is perfect, but they are running in production today, real customers are using them, and they are moving the outcomes the project set out to move.

\section{Reflections on What Could Have Been Done Better}
Hindsight makes a few things look obviously fixable. The WhatsApp automation is the clearest example -- it went through three near-complete rewrites over the course of the internship. A clearer design conversation before the first line of WhatsApp code would have removed at least one of those rewrites. The lesson I've taken away from that is simple, and Mr.~Shah had actually said it out loud at the start: for any non-trivial subsystem, a one-pager design review with the lead engineer before the first commit is cheaper than two rewrites afterwards.

The CSV upload flow for campaign creation is another. The first version accepted exactly one fixed column schema, which meant any tenant onboarding with a slightly different column order needed a code change to be supported. Generalising the input format up front is almost always cheaper than catching every variant later, especially when the data shape comes from the customer.

The desktop application is the third. The first build embedded its automation server on a hard-coded port. That worked fine on three of the four operator machines we tested on. On the fourth it collided with another running process and surfaced a cryptic error that did not actually mention a port conflict. I lost the better part of an afternoon to that one. The free-port-finding logic that the current build uses ought to have been there in version one, not bolted on as a hotfix later.

\section{Future Scope}
There is a fair bit of work the team has lined up for the next two quarters. The internal voice AI platform today runs in Indian English with code-switching into Hindi. The next phase brings dedicated agents in Tamil, Telugu, Marathi, Gujarati and Bengali. The workflow engine, currently configured per tenant statically, is moving towards dynamic branching on AA-derived signals -- a case where the customer's bank account shows a recent salary credit gets routed to a soft WhatsApp touch, while an empty account is routed straight to a human agent who can actually negotiate a settlement. A lightweight real-time sentiment classifier over the per-turn transcript would let the AI agent change tack mid-call when the customer turns hostile -- that's another piece on the roadmap. The one-off analysis on the customer records is being productionised as a predictive prioritisation model that ranks cases by their expected recovery probability. A predictive WhatsApp number-health score would let the team manage the sender pool proactively, rather than the current reactive ban-detection. A self-service lender-side allocation upload portal removes a recurring source of friction at cycle start. And a continuous voice-agent evaluation pipeline against a fixed corpus of test conversations would let the team catch prompt regressions before they reach production.

\section{Closing Note}
On a personal note: this internship has been, by some distance, the single most useful piece of engineering practice in my undergraduate degree. Getting to ship code into a system that real customers are using, being on call for the consequences of changes I had merged in, sitting through code reviews where my suggestions sometimes won and sometimes lost, and watching a quarterly review where the numbers I had contributed to were called out by name -- these are experiences a classroom course just does not give you, and there is no shortcut to them. I'm genuinely grateful to Ignosis, to my industry mentor Mr.~Rahi Shah, to my PDEU supervisor Dr.~Punit Gupta, and to the comprehensive-project programme at PDEU for making any of it possible.
